%% toc: P15 161 161.10,161.11,161.13
% Proofs for the type invariance of domains of attraction and for the stable limit of symmetric power-tail sums.
\section{Normal attraction in context: proofs}
\label{sec:annex-p15-w02c-p15-016}

\begin{annexitem}{annex:p15-w02c-p15-016-types}{thm:p15-w02c-p15-016-types}{Affine changes of the limit law}
\begin{proof}
Let \(a>0\), \(b\in\mathbb R\) and \(W(x)=V(ax+b)\). The map
\(y\mapsto ay+b\) is a continuous strictly increasing bijection of
\(\mathbb R\) onto itself. Hence \(W\) is nondecreasing and right
continuous, with limits \(0\) at \(-\infty\) and \(1\) at \(+\infty\),
so \(W\) is a distribution function. As \(y\uparrow x\), also
\(ay+b\uparrow ax+b\), so \(W(x-)=V((ax+b)-)\). Therefore \(x\) is a
continuity point of \(W\) exactly when \(ax+b\) is a continuity point
of \(V\).

Put \(B'_n=aB_n\), which is real and nonzero, and
\(A'_n=(A_n+b)/a\). Then
\[
 \zeta'_n=\frac1{B'_n}\sum_{k=1}^n\xi_k-A'_n
 =\frac{\zeta_n-b}a .
\]
Since \(a>0\), for every real \(x\),
\[
 \mathbb P(\zeta'_n\leq x)=\mathbb P(\zeta_n\leq ax+b).
\]
If \(x\) is a continuity point of \(W\), then \(ax+b\) is a continuity
point of \(V\), and the attraction of \(F\) to \(V\) gives
\(\mathbb P(\zeta_n\leq ax+b)\to V(ax+b)=W(x)\). Thus the distribution
functions of \(\zeta'_n\) converge to \(W\) at every continuity point
of \(W\), and \(F\) is attracted to \(W\) with the constants
\(A'_n,B'_n\).

If \(V\) and \(W\) are of the same type, with \(W(x)=V(ax+b)\), then
also \(V(x)=W(a^{-1}x-a^{-1}b)\) with \(a^{-1}>0\). The first part,
applied in both directions, shows that every law attracted to one of
them is attracted to the other. The two domains are therefore equal,
and the domain of attraction is the same for every law of a type.
\end{proof}
\end{annexitem}

\begin{annexitem}{annex:p15-w02c-p15-016-stable}{thm:p15-w02c-p15-016-stable}{Symmetric power tails}
\begin{proof}
Write \(X=X_1\) and let \(\rho\) be the Borel measure on
\(\mathbb R\) with density
\[
 r(u)=\frac\alpha2|u|^{-\alpha-1}\mathbf1_{\{|u|>1\}} .
\]
Since \(\int_1^\infty\alpha u^{-\alpha-1}\,du=1\) and \(r\) is even,
\(\rho\) is a symmetric probability measure, and
\(\rho((x,\infty))=\rho((-\infty,-x))=x^{-\alpha}/2\) for
\(x\geq1\).

\emph{The two forms of the hypothesis, and the law of \(X\).}
Suppose first that \(X\) and \(-X\) are equal in law and
\(\mathbb P(|X|>x)=x^{-\alpha}\) for \(x\geq1\). For \(x\geq1\), the
events \(\{X>x\}\) and \(\{X<-x\}\) are disjoint with union
\(\{|X|>x\}\), and symmetry gives
\(\mathbb P(X>x)=\mathbb P(-X>x)=\mathbb P(X<-x)\). Each therefore has
probability \(x^{-\alpha}/2\), which is (3.8.1).

Now assume only (3.8.1):
\(\mathbb P(X>x)=\mathbb P(X<-x)=x^{-\alpha}/2\) for \(x\geq1\).
We show that the distribution function \(G(y)=\mathbb P(X\leq y)\)
equals \(\rho((-\infty,y])\) for every real \(y\).
If \(y\geq1\), then \(G(y)=1-y^{-\alpha}/2\), which is the value for
\(\rho\). If \(-1\leq y<1\), then
\(\{X<-1\}\subseteq\{X\leq y\}\subseteq\{X>1\}^c\), so
\(1/2\leq G(y)\leq1-1/2\); hence \(G(y)=1/2\), again the value for
\(\rho\). If \(y<-1\), put \(x=-y>1\). The events \(\{X<-z\}\),
\(1<z<x\), decrease as \(z\uparrow x\) with intersection
\(\{X\leq-x\}\), so by continuity of probability
\(G(y)=\lim_{z\uparrow x}z^{-\alpha}/2=x^{-\alpha}/2\), which equals
\(\int_{-\infty}^y r(u)\,du=\rho((-\infty,y])\). The half-lines
\((-\infty,y]\) form a \(\pi\)-system generating the Borel sets, so
\cref{thm:p15-measure-uniqueness} shows that the law of \(X\) is
\(\rho\). In particular \(X\) is symmetric and
\(\mathbb P(|X|>x)=2\cdot x^{-\alpha}/2=x^{-\alpha}\) for \(x\geq1\).
So the two forms of the hypothesis are equivalent, and under either
the law of each \(X_j\) is \(\rho\).

\emph{Part 1.} Under \(\rho\),
\[
 \varphi(t)=\int(\cos tu+i\sin tu)\,r(u)\,du .
\]
Both integrands are bounded and \(r\) is integrable. The sine term is
an odd function times an even one, so its integral is zero. The
cosine term is even, and therefore
\[
 \varphi(t)=\alpha\int_1^\infty\cos(tu)\,u^{-\alpha-1}\,du,\qquad
 1-\varphi(t)=\alpha\int_1^\infty(1-\cos tu)\,u^{-\alpha-1}\,du,
\]
using \(\alpha\int_1^\infty u^{-\alpha-1}\,du=1\). In particular
\(\varphi\) is real, and since \(1-\cos\) is even,
\(1-\varphi(t)=1-\varphi(|t|)\).

Put \(g(u)=(1-\cos u)u^{-\alpha-1}\) for \(u>0\); it is continuous
and nonnegative. For \(s\geq0\), \(s-\sin s\) has derivative
\(1-\cos s\geq0\) and vanishes at \(0\), so \(\sin s\leq s\), and
\[
 1-\cos u=\int_0^u\sin s\,ds\leq\int_0^us\,ds=\frac{u^2}2
 \qquad(u\geq0).
\]
Hence \(g(u)\leq u^{1-\alpha}/2\) on \((0,1]\), which is integrable
there because \(1-\alpha>-1\), that is \(\alpha<2\). On
\([1,\infty)\), \(g(u)\leq2u^{-\alpha-1}\), which is integrable because
\(\alpha>0\). So \(C=\alpha\int_0^\infty g(u)\,du<\infty\). On
\((0,2\pi)\), \(\cos u<1\), so \(g>0\) there, and by continuity
\(\int_0^{2\pi}g>0\). Thus \(0<C<\infty\).

\emph{Part 2.} At \(t=0\) both sides are zero, since
\(\varphi(0)=1\). Let \(t\neq0\) and \(s>0\). In
\(1-\varphi(s)=\alpha\int_1^\infty(1-\cos sx)x^{-\alpha-1}\,dx\)
substitute \(u=sx\), so \(dx=du/s\) and
\(x^{-\alpha-1}=s^{\alpha+1}u^{-\alpha-1}\):
\[
 1-\varphi(s)=\alpha s^\alpha\int_s^\infty g(u)\,du .
\]
Take \(s=s_n=|t|n^{-1/\alpha}\); then \(s_n^\alpha=|t|^\alpha/n\), and
by the evenness from Part~1,
\[
 n\{1-\varphi(t/n^{1/\alpha})\}
 =n\{1-\varphi(s_n)\}
 =|t|^\alpha\,\alpha\int_{s_n}^\infty g(u)\,du .
\]
As \(n\to\infty\), \(s_n\downarrow0\), and the integrands
\(g\mathbf1_{(s_n,\infty)}\) increase to \(g\) on \((0,\infty)\).
Monotone convergence gives the limit \(C|t|^\alpha\).

\emph{Part 3.} For each \(k\), \(S_k\) is a Borel function of
\(X_1,\ldots,X_k\), so it is independent of \(X_{k+1}\) by
\cref{thm:p15-independent-blocks}. By
\cref{thm:p15-characteristic-algebra}, the characteristic function of
an independent sum is the product, and \(cZ\) has transform
\(t\mapsto\varphi_Z(ct)\). Induction on \(n\) gives
\[
 \mathbb E\exp(itS_n/n^{1/\alpha})=\varphi(t/n^{1/\alpha})^n
 =\Bigl(1+\frac{c_n}n\Bigr)^n,\qquad
 c_n=-n\{1-\varphi(t/n^{1/\alpha})\}.
\]
By Part~2, \(c_n\to-C|t|^\alpha\), and
\cref{thm:p15-clt-complex-products} gives, for every real \(t\),
\[
 \mathbb E\exp(itS_n/n^{1/\alpha})\to\psi(t)=\exp(-C|t|^\alpha).
\] The function \(\psi\) is continuous at \(0\),
because \(|t|^\alpha\to0\) as \(t\to0\). By the converse half of
\cref{thm:p15-characteristic-continuity-theorem}, the laws of
\(S_n/n^{1/\alpha}\) converge weakly to a unique probability law
\(\mu_Y\) with characteristic function \(\psi\). Let \(Y\) be the
identity map on \((\mathbb R,\mathcal B(\mathbb R),\mu_Y)\); then
\(S_n/n^{1/\alpha}\Rightarrow Y\) and \(\mathbb Ee^{itY}=\psi(t)\).

\emph{Part 4.} By \cref{thm:p15-characteristic-normal-transform}, a
normal law with real mean \(m\) and variance \(v\geq0\) has transform
\(\exp(imt-vt^2/2)\), of modulus \(\exp(-vt^2/2)\); the case \(v=0\)
is the point mass at \(m\), of modulus one. A law determines its
characteristic function, so if the law of \(Y\) were \(N(m,v)\), then
\[
 \exp(-C|t|^\alpha)=\exp(-vt^2/2)\qquad(t\in\mathbb R).
\]
At \(t=1\) this gives \(C=v/2\), so \(v=2C>0\). At \(t=2\) it gives
\(2^\alpha C=2v=4C\), so \(2^\alpha=4\) and \(\alpha=2\), contradicting
\(\alpha<2\). Hence the law of \(Y\) is no normal law. Taking
\(v=0\) shows that it is not a point mass, so it is nondegenerate.
Taking \(m=0\), \(v=1\) shows that it is not the law of \(\chi\).
\end{proof}
\end{annexitem}

\begin{annexitem}{annex:p15-w02c-p15-016-cauchy-case}{thm:p15-w02c-p15-016-cauchy-case}{The constant and the limit at index one}
\begin{proof}
Take \(\alpha=1\), so \(C=\int_0^\infty(1-\cos u)u^{-2}\,du\). By
\cref{thm:p15-characteristic-inverse-examples}, the even function
\(f_P(x)=(1-\cos x)/(\pi x^2)\), \(f_P(0)=1/(2\pi)\), is the density of
a law with characteristic function \((1-|t|)^+\). Evaluating that
transform at \(t=0\) gives \(\int_{\mathbb R}f_P=1\). Evenness gives
\(\int_0^\infty f_P=1/2\), so
\(C=\pi\int_0^\infty f_P(u)\,du=\pi/2\).

By \cref{thm:p15-w02c-p15-016-stable},
\(S_n/n\Rightarrow Y\) with \(\mathbb Ee^{itY}=e^{-\pi|t|/2}\). The
standard Cauchy variable \(Z\) has transform \(e^{-|t|}\) by
\cref{thm:p15-characteristic-inverse-examples}, so by
\cref{thm:p15-characteristic-algebra} \((\pi/2)Z\) has transform
\(e^{-|\pi t/2|}=e^{-\pi|t|/2}\). By
\cref{thm:p15-characteristic-inversion-uniqueness}, \(Y\) and
\((\pi/2)Z\) have the same law, which proves the displayed limit.

Under the law \(\rho\) of the preceding proof with \(\alpha=1\),
\[
 \mathbb E|X_1|=2\int_1^\infty u\cdot\frac12u^{-2}\,du
 =\int_1^\infty\frac{du}u=\infty .
\]
Finally, suppose \(S_n/n\to c\) in probability for a constant \(c\).
For \(t\in\mathbb R\) and \(\varepsilon>0\), since
\(|e^{iy}-1|\leq\min(|y|,2)\),
\[
 \bigl|\mathbb Ee^{itS_n/n}-e^{itc}\bigr|
 \leq\mathbb E\bigl|e^{it(S_n/n-c)}-1\bigr|
 \leq|t|\varepsilon+2\,\mathbb P(|S_n/n-c|>\varepsilon),
\]
so \(\mathbb Ee^{itS_n/n}\to e^{itc}\), which has modulus one. The
limit \(e^{-\pi|t|/2}\) has modulus \(e^{-\pi/2}<1\) at \(t=1\), so no
such constant exists.
\end{proof}
\end{annexitem}
